% ============================================================
%  6.4 Métricas de evaluación   (capítulo 6, Marco teórico)
%  Se incluye desde secciones/07_marco_teorico.tex con \input.
% ============================================================
\section{Métricas de evaluación}
\label{sec:metricas}

\subsection{Modelo de referencia}
\label{subsec:modelo-referencia}

Toda evaluación de modelos requiere un punto de comparación, porque ninguna métrica de error
tiene sentido en términos absolutos: una devianza o un error medio solo son grandes o
pequeños en relación con lo que lograría una alternativa razonable, y dependen de la escala
de los conteos, del número de zonas y de la dispersión intrínseca de los datos. En el
proyecto, el modelo de referencia es la tasa global proporcional a la población descrita en
\ref{subsec:lineas-base}, $\hat{\mu}_i = \hat{\lambda} \cdot P_i$. Su elección no es
arbitraria, sino que materializa la pregunta del proyecto: si un modelo no mejora esta
referencia, el contexto territorial no aporta información sobre las consultas esperadas más
allá de la población, y la brecha observado-esperado se reduce a una comparación de tasas
brutas. Todas las métricas se reportan, por tanto, junto con su valor para el modelo de
referencia y, cuando procede, como mejora relativa respecto de él.

\subsection{Devianza de Poisson}
\label{subsec:devianza-poisson}

La devianza de Poisson mide la discrepancia entre los conteos observados y los esperados en
la escala natural de los datos de conteo. Se define como
\[
  D = 2 \sum_i \left[ y_i \ln\!\left(\frac{y_i}{\hat{\mu}_i}\right) - (y_i - \hat{\mu}_i) \right],
\]
con la convención de que $y_i \ln(y_i / \hat{\mu}_i) = 0$ cuando $y_i = 0$, y la devianza
media como $D / n$. Equivale al doble de la diferencia entre la log-verosimilitud del modelo
saturado y la del modelo evaluado, por lo que constituye la medida de ajuste natural de los
modelos de Poisson \parencite{mccullagh1989}. En scikit-learn, la función
\var{mean_poisson_deviance} calcula la devianza media y equivale a la devianza de Tweedie
con parámetro de potencia igual a uno; requiere $y \ge 0$ y predicciones estrictamente
positivas \parencite{sklearn-mpd}.

La preferencia por la devianza de Poisson frente al error absoluto medio (MAE) o a la raíz
del error cuadrático medio (RMSE) obedece a la heterogeneidad de escala de los conteos. En un
conjunto de zonas con conteos que van desde cero hasta valores elevados, el RMSE queda
dominado por los errores de las pocas zonas con muchas consultas, mientras que la devianza
de Poisson pondera cada error en relación con la magnitud esperada, en coherencia con la
propiedad de que la varianza de un conteo crece con su media. Además, la devianza de Poisson
es una función de pérdida estrictamente consistente para la media, en el sentido de
\textcite{gneiting2011}: su valor esperado se minimiza cuando la predicción coincide con el
verdadero valor esperado, que es precisamente el estimando del proyecto. El MAE, en cambio,
es consistente para la mediana, que en conteos pequeños y asimétricos puede diferir
notablemente de la media.

\subsection{Fracción de devianza explicada (\texorpdfstring{$D^2$}{D²})}
\label{subsec:d2}

La fracción de devianza explicada, o $D^2$, es el análogo del coeficiente de determinación
$R^2$ para modelos de conteo, y se define como
$D^2 = 1 - D_{\mathrm{modelo}} / D_{\mathrm{nulo}}$, donde $D_{\mathrm{nulo}}$ es la devianza
de un modelo nulo de comparación. Un valor de uno indica ajuste perfecto, un valor de cero
indica que el modelo no mejora al nulo, y un valor negativo indica que es peor. En
scikit-learn, la función \var{d2_tweedie_score} con parámetro de potencia igual a uno calcula
esta fracción tomando como modelo nulo la media constante de los valores observados; como en
el caso del $R^2$, la puntuación puede ser negativa \parencite{sklearn-d2}.

La interpretación del $D^2$ exige declarar contra qué modelo nulo se compara, porque su valor
cambia radicalmente según la elección. Frente a una media constante, que asigna el mismo
número de consultas a todas las zonas con independencia de su población, cualquier modelo
que incorpore la población obtendrá un $D^2$ elevado, sin que ello revele valor alguno del
contexto territorial. Frente al modelo de referencia con población, el $D^2$ mide la mejora
atribuible al contexto. Por ello, en el proyecto se reportan ambos valores de forma
explícita y se añade la mejora relativa frente al modelo de referencia,
$\mathrm{MR} = 1 - D_{\mathrm{modelo}} / D_{\mathrm{referencia}}$, que expresa qué fracción
de la devianza no explicada por la tasa global logra reducir cada técnica.

\subsection{Calibración}
\label{subsec:calibracion}

Un modelo puede ordenar bien las zonas y, sin embargo, sobrestimar o subestimar
sistemáticamente el nivel de las consultas. La calibración global se evalúa mediante el
cociente $C = \sum_i \hat{\mu}_i / \sum_i y_i$, que debe ser cercano a uno; los modelos de
Poisson con intercepto ajustados por máxima verosimilitud satisfacen esta igualdad en los
datos de entrenamiento, pero no necesariamente en datos no vistos. La calibración por
deciles complementa esta medida: se ordenan las zonas según su valor predicho, se agrupan en
diez grupos de igual tamaño y se compara, en cada grupo, la media de las consultas predichas
con la media de las observadas. La representación gráfica de estos pares frente a la
diagonal permite detectar si el modelo sobrestima las zonas de baja demanda esperada o
subestima las de alta demanda, patrones que afectarían directamente a la ordenación de las
brechas.

\subsection{Residuos para conteos}
\label{subsec:residuos}

El análisis de residuos permite diagnosticar dónde y cómo falla un modelo. El residuo de
Pearson se define como $r_{P,i} = (y_i - \hat{\mu}_i) / \sqrt{V(\hat{\mu}_i)}$, donde
$V(\mu)$ es la función de varianza del modelo, igual a $\mu$ en la Poisson y a
$\mu + \alpha\mu^2$ en la binomial negativa; la suma de sus cuadrados, dividida por los
grados de libertad residuales, proporciona el estadístico de dispersión de
\ref{subsec:datos-conteo}. El residuo de devianza se define como
$r_{D,i} = \operatorname{signo}(y_i - \hat{\mu}_i) \cdot \sqrt{d_i}$, con
$d_i = 2\left[y_i \ln(y_i / \hat{\mu}_i) - (y_i - \hat{\mu}_i)\right]$, de modo que la suma
de sus cuadrados es la devianza total \parencite{mccullagh1989}. Los residuos de devianza
tienen, en general, una distribución más próxima a la simétrica que los de Pearson, lo que
facilita su inspección visual.

En un contexto espacial, los residuos deben examinarse también en el mapa. La
representación cartográfica de los residuos de las zonas de validación permite identificar
áreas en las que el modelo sobrestima o subestima de forma sistemática, y el cálculo del $I$
de Moran y de los indicadores LISA sobre los residuos permite diagnosticar si persiste
autocorrelación espacial no explicada \parencite{anselin1995}. Una autocorrelación residual
significativa indica que falta en el modelo alguna variable con estructura espacial y, en el
contexto del proyecto, puede señalar agrupamientos de zonas con déficit de consultas
compatibles con áreas enteras de cartografía incompleta, información sustantiva para la
priorización.

\subsection{Estabilidad de la priorización}
\label{subsec:estabilidad-priorizacion}

Dado que el producto final del proyecto es una ordenación de zonas para su revisión, importa
tanto la exactitud de las estimaciones como la estabilidad de esa ordenación frente a
decisiones metodológicas razonables. La correlación de rangos de Spearman
\parencite{spearman1904} mide la concordancia entre dos ordenaciones, y en ausencia de
empates se calcula como
\[
  \rho = 1 - \frac{6 \sum_i d_i^2}{n(n^2 - 1)},
\]
donde $d_i$ es la diferencia entre los rangos de la zona~$i$ en ambas ordenaciones y $n$ el
número de zonas. Un valor cercano a uno indica que las dos variantes ordenan las zonas de
forma casi idéntica.

Como la priorización se concentra en el extremo superior de la lista, la correlación de
rangos, que pondera por igual todo el ordenamiento, se complementa con una medida de
solapamiento del conjunto de zonas priorizadas. Si $A$ y $B$ son los conjuntos de las $k$
zonas con mayor prioridad bajo dos variantes, el índice de Jaccard
$J = \lvert A \cap B \rvert / \lvert A \cup B \rvert$ mide la proporción de zonas
compartidas. En el proyecto, estas medidas se calculan entre las variantes que surgen de las
decisiones metodológicas documentadas: el período anterior al hueco de siete semanas, un
umbral de población mínima más exigente y la técnica de estimación. Las variantes con las
dos ediciones de Kontur Population (2022-06-30 y 2023-11-01) y con las resoluciones H3 7 y
9, comparadas tras agregar las prioridades a una escala común, se recomiendan como
extensión.

\subsection{Comparación de grupos con y sin cobertura}
\label{subsec:comparacion-cobertura}

Para describir si las zonas cubiertas por el trazado de las rutas difieren de las no
cubiertas en su tasa de consultas o en su brecha, el punto de partida es la prueba de
Mann-Whitney \parencite{mann1947}, una prueba no paramétrica que no supone normalidad y que contrasta si
los valores de un grupo tienden a ser mayores que los del otro. Su estadístico $U$ cuenta el
número de pares, formados por una observación de cada grupo, en los que la del primer grupo
supera a la del segundo, y resulta adecuado para distribuciones asimétricas con abundancia de
ceros como las de este proyecto.

La significación estadística se complementa con un tamaño de efecto. El $\delta$ de Cliff
\parencite{cliff1993} se define como
\[
  \delta = \frac{\#(x_i > y_j) - \#(x_i < y_j)}{n_1 \cdot n_2},
\]
es decir, la diferencia entre la proporción de pares en los que el primer grupo supera al
segundo y la proporción en que ocurre lo contrario; toma valores entre $-1$ y $1$ y se
relaciona con el estadístico de Mann-Whitney mediante $\delta = 2U / (n_1 n_2) - 1$. Para su
interpretación se emplean los umbrales de \textcite{romano2006}: efecto despreciable si
$\lvert\delta\rvert < 0{,}147$, pequeño si $\lvert\delta\rvert < 0{,}33$, mediano si
$\lvert\delta\rvert < 0{,}474$ y grande en caso contrario. Debe advertirse que, con
autocorrelación espacial, las observaciones no son independientes y el tamaño efectivo de la
muestra es menor que el nominal, de modo que los $p$-valores de la prueba de Mann-Whitney
resultan optimistas; por ello, en el proyecto se reporta solo la magnitud del efecto
($\delta$ de Cliff), sin $p$-valor, y la comparación se interpreta como descriptiva, no como
prueba causal del efecto de la cobertura.

\subsection{Evaluación desagregada y análisis de sensibilidad}
\label{subsec:evaluacion-desagregada}

Una métrica global puede ocultar desempeños muy desiguales entre subconjuntos del
territorio. Por ello, las métricas de \ref{subsec:devianza-poisson} a
\ref{subsec:calibracion} se calculan también de forma desagregada: por anillos de distancia
al centro, para verificar si el modelo funciona igual en el centro que en la periferia; por
municipio, para detectar si alguno queda sistemáticamente mal representado; y por niveles de
población (por ejemplo, terciles), para examinar el comportamiento en las zonas menos
pobladas, donde el problema de los números pequeños es más agudo. Un modelo cuya ventaja
global se concentre en un solo estrato, o que empeore en las zonas periféricas donde la
priorización es más necesaria, debe interpretarse con cautela.

El análisis de sensibilidad completa la evaluación examinando la robustez de las
conclusiones frente a las decisiones metodológicas que no pueden fijarse con certeza. En
este ciclo se evalúan el período de datos respecto del hueco de siete semanas, el umbral de
población mínima y la técnica de estimación. La edición de la población, la resolución de
la grilla, el umbral de cobertura y el tamaño de los bloques de validación se recomiendan
como extensión. Su propósito no es elegir la variante más favorable, sino documentar qué
conclusiones se mantienen bajo todas las variantes razonables y cuáles dependen de supuestos
particulares, información indispensable para que los voluntarios interpreten la lista de
zonas prioritarias con el grado de confianza adecuado.
